%============================================================
% Chapter 13: Tree Fundamentals
% STATUS: PLANNED  (generated by tools/outline.py; replace this file when the chapter is written)
%============================================================
\chapter{Tree Fundamentals}
\label{ch:tree_basics}

\begin{tcolorbox}[enhanced, colback=white, colframe=secondaryblue, boxrule=1pt, arc=2mm]
\textbf{\large\textcolor{primaryblue}{Trees are the graphs you will meet most: root them, measure them, and recognise them up to isomorphism.}}
\end{tcolorbox}

\chapterfacts{1400}{2200}{5}{10}{30}

\begin{plannedbox}[Planned Contents --- this chapter is not written yet]
\textbf{Part A --- Lessons}
\begin{enumerate}[leftmargin=*, label=\textbf{A\arabic*.}, itemsep=1pt]
    \item Rooting a tree: parents, depths, subtree sizes
    \item Diameter: two searches, and the DP view
    \item Centres and centroids
    \item The Euler tour: flattening a tree so every subtree is a contiguous range (ancestor tests were Chapter 2; range queries are Chapter 15)
    \item Tree isomorphism by hashing; Pr\"ufer sequences and counting labelled trees
\end{enumerate}
\textbf{Part B --- Patterns} (with candidate problems, checked when the chapter is written)
\begin{enumerate}[leftmargin=*, label=\textbf{P\arabic*.}, itemsep=1pt]
    \item \textbf{Diameter and eccentricity} --- Tree Diameter {\small\color{darkgray}(CSES 1131)}, Tree Distances I {\small\color{darkgray}(CSES 1132)}
    \item \textbf{Centres and centroids} --- Link Cut Centroids {\small\color{darkgray}(CF 1406C)}
    \item \textbf{Euler tour flattening}
    \item \textbf{Isomorphism and hashing} --- Symmetree {\small\color{darkgray}(CF 1800G)}
    \item \textbf{Counting trees} --- Clues {\small\color{darkgray}(CF 156D)}
\end{enumerate}
\textbf{Connections.} Tree DP is \dpbook{11}{DP on Trees}; rerooting is \dpbook{19}{Rerooting DP}.
\end{plannedbox}
